\documentclass[11pt]{article}
\usepackage[margin=1in]{geometry}
\usepackage{amsmath,amssymb,amsthm,mathtools}
\usepackage[hidelinks]{hyperref}

\newtheorem{theorem}{Theorem}[section]
\newtheorem{proposition}[theorem]{Proposition}
\newcommand{\splus}{s^+}
\newcommand{\sminus}{s^-}
\newcommand{\sig}{\sigma}
\newcommand{\one}{\mathbf 1}

\title{Positive Square Energy of Every Cyclomatic-Rank-Twelve Cactus}
\author{AI-generated manuscript; no human author is claimed}
\date{30 July 2026}

\begin{document}
\maketitle

\begin{abstract}
We prove that every connected cactus $G$ of cyclomatic rank twelve on $n$
vertices satisfies $s^+(G)>n$.  The sharp cactus doubly nonnegative formula
leaves exactly the cycle multisets $T^{11}Q$ and $T^{10}PP$, where
$T=C_3$, $P=C_5$, and $Q$ is an arbitrary cycle.  The all-rank hostile and
nonhostile one-cycle theorems close the first family for every parity of $Q$;
the all-rank two-pentagon theorem closes the second.  Only the exact frontier
calculation is repeated here; the three rank-uniform structural proofs are
cited rather than duplicated.
\end{abstract}

\section{Sharp DNN bound and exact frontier}

All graphs are finite, simple, and undirected.  A cactus is connected and has
only bridge or cycle blocks; its cyclomatic rank is the number of cycle blocks.
For adjacency eigenvalues $\lambda_1,\ldots,\lambda_n$, put
\[
 \splus(G)=\sum_{\lambda_i>0}\lambda_i^2,
 \qquad
 \sminus(G)=\sum_{\lambda_i<0}\lambda_i^2,
 \qquad
 \sig(G)=\splus(G)-|V(G)|.
\]

For an edge-containing graph define
\[
 \kappa(G)=\sup_{\substack{M\succeq0,\ M\geq0\\\one^TM\one>0}}
 \frac{4\bigl(\sum_{uv\in E(G)}\sqrt{M_{uv}}\bigr)^2}
      {\one^TM\one}.
\]
The sharp cactus DNN theorem~\cite{DNN,LTZ} states that if a cactus has $b$
bridge blocks and cycle blocks $C_{\ell_1},\ldots,C_{\ell_r}$, then
\begin{equation}\label{eq:kappa}
 \sminus(G)\leq\kappa(G)
 =b+\sum_{i=1}^r\kappa(C_{\ell_i}),
 \qquad
 \kappa(C_\ell)=
 \begin{cases}
  \ell,&\ell\text{ even},\\
  \ell\sec^2\!\bigl(\pi/(2\ell)\bigr),&\ell\text{ odd}.
 \end{cases}
\end{equation}
Set
\[
 \epsilon_\ell=
 \begin{cases}
  0,&\ell\text{ even},\\
  \ell\tan^2\!\bigl(\pi/(2\ell)\bigr),&\ell\text{ odd}.
 \end{cases}
\]
For rank twelve, $|E(G)|=n+11$ and block counting gives
$b+\sum_i\ell_i=n+11$.  Since
$\splus(G)+\sminus(G)=2|E(G)|$, equation~\eqref{eq:kappa} yields the sharp-DNN
surplus formula
\begin{equation}\label{eq:surplus}
 \boxed{\sig(G)\geq11-\sum_{i=1}^{12}\epsilon_{\ell_i}.}
\end{equation}

\begin{proposition}[Exact rank-twelve frontier]\label{prop:frontier}
The right side of~\eqref{eq:surplus} is nonpositive exactly for
\[
 \boxed{T^{11}Q}\quad(q\geq3,\text{ including }Q=T),
 \qquad
 \boxed{T^{10}PP}.
\]
Every other rank-twelve cycle multiset has positive surplus.
\end{proposition}

\begin{proof}
For real $x\geq3$, $\epsilon_x=x\tan^2(\pi/(2x))$ decreases strictly: with
$u=\pi/(2x)$, the derivative of $\tan^2u/u$ has the sign of
$2u-\sin u\cos u>0$.  Thus every nontriangle contributes at most
\[
 a:=\epsilon_5=5-2\sqrt5,
\]
and the exact comparisons needed below are
\begin{equation}\label{eq:exact}
 0<a<1,\qquad3a<2,\qquad2a>1,
 \qquad a+\epsilon_7<1.
\end{equation}
The first four inequalities after moving radicals to positive sides have
squaring certificates $20<25$, $4<5$, $169<180$, and $80<81$.  For the last,
$\cos(\pi/7)>2159/2401>7/8$ gives $\epsilon_7<7/15$, while
$7/15<2\sqrt5-4=1-a$ follows from $4489<4500$.

Let $t$ be the number of triangles.  If $t\leq9$, monotonicity and
\eqref{eq:exact} give
$\sum_i\epsilon_{\ell_i}\leq9+3a<11$.  If $t=10$, the two nontriangles
contribute less than one when one is even, or when their odd lengths are not
both five, since then their sum is at most $a+\epsilon_7<1$.  The unique
remaining pair is $PP$, and $10+2a>11$.  If $t\geq11$, the multiset is
$T^{11}Q$ and its excess sum is $11+\epsilon_q\geq11$.  These cases are
exhaustive.  The exact note and fail-closed arithmetic verifier independently
freeze this classification~\cite{Frontier}.
\end{proof}

\section{Closure of the frontier}

\begin{proposition}\label{prop:onecycle}
Every connected cactus whose cycle blocks are eleven triangles and one cycle
$Q=C_q$ satisfies $\splus(G)>|V(G)|$.
\end{proposition}

\begin{proof}
If $q\equiv1\pmod4$, then $q\geq5$ and the all-rank
triangle--hostile-cycle theorem applies~\cite{Hostile}.  If $q$ is even or
$q\equiv3\pmod4$, the all-rank nonhostile one-cycle theorem applies
\cite{Nonhostile}.  These alternatives exhaust all integers $q\geq3$, and the
nonhostile theorem explicitly includes $q=3$.  Both theorems allow arbitrary
shared cuts, bridge connectors, and finite tree attachments.
\end{proof}

\begin{proposition}\label{prop:twopent}
Every connected cactus whose cycle blocks are ten triangles and two pentagons
satisfies $\splus(G)>|V(G)|$.
\end{proposition}

\begin{proof}
This is the all-rank two-pentagon theorem~\cite{TwoPent}, specialized to
$r=10$.  Its statement permits arbitrary cactus incidence, bridge lengths and
branching, and arbitrary finite tree attachments.
\end{proof}

\begin{theorem}\label{thm:main}
Every connected cyclomatic-rank-twelve cactus $G$ on $n$ vertices satisfies
\[
 \boxed{\splus(G)>n}.
\]
\end{theorem}

\begin{proof}
Outside the two families in Proposition~\ref{prop:frontier}, the strict
inequality follows immediately from~\eqref{eq:surplus}.  Proposition
\ref{prop:onecycle} closes $T^{11}Q$ for every parity of $Q$, and Proposition
\ref{prop:twopent} closes $T^{10}PP$.  The frontier classification is
exhaustive, so the theorem follows.
\end{proof}

\section{Reproduction, scope, and disclosure}

From the public repository root, reproduce the exact frontier audit with
Python 3.10 or newer:
\begin{verbatim}
python3 research/rank-twelve-cactus-dnn-residual-frontier-verifier.py
python3 -O research/rank-twelve-cactus-dnn-residual-frontier-verifier.py
\end{verbatim}
Both runs report the certificate SHA-256 digest
\begin{center}
\texttt{96d20340187cd0b2c01ca5d89d1f2c06cb0ecd321d035d73b53d94a706edd663}.
\end{center}
The SHA-256 digest of the cited verifier file itself is
\begin{center}
\texttt{d94364685251afb9a1045085f6cd4b85774d37329bbf359a1c35dd738c3baae6}.
\end{center}
The verifier certifies only the sharp-DNN frontier; positivity on the two
residual families comes from the cited all-rank manuscripts.

\paragraph{Scoped nonclaim.}
This paper proves only the stated result for finite simple connected cacti of
cyclomatic rank twelve.  It does not claim the positive-square-energy
conjecture for all graphs, nor a theorem for arbitrary cactus rank or for
non-cactus block intersections.  The cited two-pentagon theorem concerns
exactly triangles and two pentagons; it makes no claim for three pentagons.

\section*{AI disclosure}
This manuscript was generated and assembled by an AI coding and mathematical
reasoning system from the cited public-repository manuscripts, note, and exact
verifier.  No human authorship, human review, or human verification is claimed.
The result should be assessed from the displayed reduction and the cited proof
objects, not from asserted authority.

\begin{thebibliography}{99}
\bibitem{LTZ} Y.~Liu, Q.~Tang, and S.~Zhang,
\emph{The positive and negative square-energy conjecture}, arXiv:2607.18031,
2026.

\bibitem{DNN} Companion manuscript, \emph{The Optimized Liu--Tang--Zhang
Constant for Cactus Graphs}, 2026; \path{sharp-cactus-dnn/paper.tex}.

\bibitem{Frontier} \emph{Exact rank-twelve cactus DNN residual frontier},
30 July 2026; \path{research/rank-twelve-cactus-dnn-residual-frontier-2026-07-30.md}
and \path{research/rank-twelve-cactus-dnn-residual-frontier-verifier.py}.

\bibitem{Hostile} Companion manuscript, \emph{Cacti with Arbitrarily Many
Triangles and One Hostile Cycle}, 2026;
\path{all-rank-triangle-hostile-cacti/paper.tex}.

\bibitem{Nonhostile} \emph{All-rank nonhostile one-cycle cactus theorem},
30 July 2026;
\path{research/all-rank-nonhostile-one-cycle-theorem-2026-07-30.md}.

\bibitem{TwoPent} Companion manuscript, \emph{Positive Square Energy for Cacti
with Triangles and Two Pentagons}, 2026;
\path{all-rank-triangle-two-pentagon-cacti/paper.tex}.
\end{thebibliography}

\end{document}
